\scanpage{86}
\mspage{\handwritten{63}}

(\textsuperscript{12}!) This is nonsense, as one sees already in the following picture where $f\ f''=\mathrm{S}_{\unsure{o|0}}$ :

(\textsuperscript{13}!) This is nonsense again for $n\quad 1$, see PS at the end of section 13.

(\textsuperscript{14}!) Same mistake as \deleted{in} the one noticed in the previous note. The fibered products exist in $\underline{\mathrm{C}}_{o}$ only, and \uline{these} should be preserved by the functors under consideration. Thus the ``universal problem'' \deleted{\unsure{(in}}\unsure{'}has to be rephrased somewhat ...

(\textsuperscript{15}) I guess they are\unsure{.} not equivalent, even when restricting to objects which are groupoids (i.e. so-called Gr-categories).

(\textsuperscript{16}) An object of $\mathrm{M}$ is called fibrant (resp. trivial) it the map from it to the final object is fibrant (resp. a weak equivalence).

(\textsuperscript{17}) This seems doubtful, unless all objects of $\mathrm{M}$ are fibrant (which is true for the most familiar cases I was having in mind). But even assuming this, we still need a unicity statement (in a suitable sense) for the functors $\underline{\mathrm{B}}_{o}\quad \mathrm{M}$ thus obtained, in order to be sure that the correspon\deleted{d}ding functors $\mathrm{M}\quad(\mathrm{Ho})$ are all canonically isomorphic. These kind of questions may be viewed as closely related to the question of existence (and unicity) of ``test functors'' for a given test category $\mathrm{A}$ (here, the category of ``standard hemispheres'' ) into a given asphericity or contractibility structure, as discussed in section 90 (without yet getting there \deleted{\unsure{the}}\inserted{any} clear-cut handy existence and unicicity theorems which are to be \deleted{\unsure{hped}} hoped for). The set-up of asphericity and contractibility structures, which has been worked out in the months following, has been gradually replacing Quillen's approach to homotopy models. (I'll have to come back in due course to theq\deleted{\unsure{z}}uestion of the relationship between the two approaches, which deserves to be understood.)

(\textsuperscript{18}) The word ``gerb'' here stands \deleted{for} as a ``translation'' of the F\shortstack{\scriptsize\unsure{''}\\r\\[-.4ex]\scriptsize\unsure{\textasciitilde}}ench word ``gerbe'', as used in Giraud's book on non commutative cohomological algebra. With the terminology of next section, we would call it rather a ``stack of groupoids'' or ``Gr-stack'' (with a specification by $n$ or $\infty$ if needed !).

(\textsuperscript{19}\deleted{\unsure{)}}\unsure{!}) This is false, as seen below.

(\textsuperscript{20}\deleted{)}!) This descr\deleted{\unsure{p}}iption is dubious, as it may give categories $\underline{\mathrm{C}}_{n}$ which are larger than the ones we want.

(\textsuperscript{21}!) It seems dubious however that the mere category structure of $\underline{\mathrm{C}}$ will allow us to recover the \deleted{\unsure{su}} ``primitive'' subcategory $\underline{\mathrm{C}}_{o}$ , and it looks safer to add the latter as an extra structure to $\underline{\mathrm{C}}$ .
% archive sheet number: 85
